\par\Needspace{7\baselineskip}\bigskip
\begin{center}
\pdfbookmark[2]{Gauss an Bessel, Göttingen, 21. März 1843}{pallas-brief-18430321}
\textsc{Gauss} an \textsc{Bessel}, Göttingen, 21. März 1843.
\end{center}
\noindent\hspace*{\parindent}\mbox{. . . . .}\quad Die erste [\textsc{Bessel}sche] Abhandlung über die Jupitersmasse hat bei mir eine Erinnerung geweckt, die mir immer schmerzhaft ist, nemlich an meine alte Arbeit über die Pallasstörungen. Sie ist seit fast einem Vierteljahrhundert mir so fremd geworden, dass es mir schwer wird, mich selbst in den vorhandenen Papieren zu orientiren. Die Nothwendigkeit einer Vergrösserung der \textsc{Laplace}schen Jupitersmasse wurde freilich sogleich erkannt, ich setzte sie anfangs auf \(\frac{48}{47}\), ohne dass ich in diesem Augenblick ausfinden kann, wann ich diese Bestimmung gemacht habe; ich finde aber später, dass aus den ersten 10 Oppositionen, \(-2376\) Tage . . . . . \(+1759{,}5\) Tage (von Anfang 1810 als Epoche gerechnet), die Elemente von Grund aus verbessert sind, und darin u.\,a. auch jener Bruch auf \(\frac{43}{42}\) vergrössert ist. Es fand sich die Summe der Quadrate der übrig bleibenden Fehler (NB. ohne Berücksichtigung anderer Störungen, namentlich der durch Saturn):
\begin{center}
\renewcommand{\arraystretch}{1.3}
\begin{tabular}{@{}lrr@{}}
\textsc{Laplaces} Masse \(\left[\frac{1}{1067{,}09}\right]\) & (Hel.) Länge \(3175{,}69\) & (Hel.) Breite \(348{,}54\) \\
vergrössert um \(\frac{1}{47}\) & \(963{,}78\) & \(318{,}47\) \\
vergrössert um \(\frac{1}{42}\) & \(918{,}46\) & \(322{,}66\).
\end{tabular}
\end{center}
\par\SourceLeaf{444}Dann findet sich noch eine neue Verbesserungsrechnung aller Elemente gestützt auf die ersten 12 Oppositionen, allein in den Papieren, die ich jetzt auffinden kann, ist die Rechnung nicht ganz zu Ende geführt. Eben bei diesem Suchen fand ich aber noch einige ältere Papiere, worauf eine ähnliche Rechnung aus den ersten 9 Oppositionen steht\HistoricalNoteEight{S. 561 dieses Bandes.]}, ohne dass ich darum mit Gewissheit sagen kann, ob nicht noch früher die Correction der Masse gefunden ist. Zugleich aber sehe ich, dass, was oben aus dem Gedächtniss geschrieben war, falsch gewesen ist, nemlich die Störungen durch Saturn sind allerdings mit berücksichtigt. Es ist wohl überflüssig zu bemerken, dass allen diesen Rechnungen diejenige Störungstheorie zum Grunde liegt, die ich die allgemeine nenne, indem danach ihr Betrag für jede beliebige Zeit gefunden wird; es sind dabei zusammen über 800 Gleichungen (wenn man diejenigen, welche Sinus und Cosinus desselben Winkels enthalten, zusammenfasst, ohne diese Zusammenfassung also über 1600), und diese Theorie war mehr als einmal ganz durchgearbeitet, was nöthig war, um erst die mittlern Elemente selbst zu erhalten. Neben dieser generellen Theorie war aber zugleich fortwährend die Rechnung durch Quadraturen fortgeführt, welche für jene eine ganz unentbehrliche Controlle lieferte. --- Das letzte, was sich vorfindet, ist die Berechnung und Vergleichung der 14\textsuperscript{ten} Opposition vom 6. Januar 1820. Ob oder wann ich dazu kommen werde, in diese Arbeiten mich wieder hineinzustudiren und sie zu redigiren und zu publiciren, steht dahin; ohne eine Abänderung meiner äussern Lage, so dass ich völlig frei über meine Zeit schalten kann, wird es schwerlich geschehen, zumal da so viele andere Dinge noch sich präsentiren, die ich eben so ungern untergehen lassen mag, oder richtiger, auf die ich selbst grössern Werth lege.
\par\bigskip\centerline{\rule{3cm}{.4pt}}\par
